\chapter{Mười Mẫu Truyện Về Học Thêm, Thành Tích Và Sự Mệt Mỏi Được Khen}
\markboth{Mười Mẫu Truyện Về Học Thêm, Thành Tích Và Sự Mệt Mỏi Được Khen}{Mười Mẫu Truyện Về Học Thêm, Thành Tích Và Sự Mệt Mỏi Được Khen}

\begin{center}
	\textit{\color{bookprimary}``Có những đứa trẻ được khen là cố gắng, trong khi thứ đang diễn ra thật ra là kiệt sức có tổ chức.''}

	\textit{\color{bookprimary}``Some children are praised as hardworking when what is really happening is organized exhaustion.''}
\end{center}

\ngancach

\latcat{1. Đứa Trẻ Ngủ Gục Trên Xe}{Tan lớp thêm lúc tối muộn, em ngồi sau xe và ngủ gật ngay khi bánh xe vừa lăn. Người lớn nhìn thấy sự vất vả đó rồi bảo: ``Ráng chút, sau này mới hơn người.'' Câu nói nghe giống động viên, nhưng cũng có thể là cách hợp thức hóa việc một đứa trẻ đang phải lớn lên trong mệt mỏi kéo dài.}{After late tutoring, the child falls asleep on the back of the motorbike as soon as the wheels move. Adults see the exhaustion and say, ``Endure it now so you can surpass others later.'' It sounds encouraging, but it can also legitimize a childhood stretched by chronic fatigue.}

\latcat{2. Cặp Sách Đầy, Đầu Óc Mỏng}{Em đi đủ lớp, mang đủ tài liệu, đổi đủ giáo viên, nhưng vẫn không giải thích nổi một ý cơ bản bằng lời của mình. Vấn đề không còn là thiếu lớp. Vấn đề là học quá nhiều bằng cách vay mượn mà quá ít bằng cách tiêu hóa.}{The child attends many classes, carries many materials, changes many teachers, yet still cannot explain a basic idea in their own words. The problem is no longer a lack of classes. The problem is learning too much by borrowing and too little by digesting.}

\latcat{3. Tự Hào Vì Không Có Ngày Nghỉ}{Có học sinh nói bằng giọng hãnh diện rằng cả tuần em không có buổi nào rảnh. Nghe rất ghê gớm. Nhưng một đứa trẻ mất sạch ngày nghỉ chưa chắc đang chiến thắng. Có khi em chỉ đang quen coi sự cạn kiệt là tiêu chuẩn đáng tự hào.}{Some students say proudly that they have no free session all week. It sounds impressive. But a child with no rest days is not necessarily winning. Sometimes the child is only being trained to mistake depletion for achievement.}

\latcat{4. Đi Học Thêm Để Khỏi Bị So}{Có em xin đi học thêm không phải vì muốn giỏi hơn mà vì sợ bị thua bạn bè. Khi động lực học bị buộc quá chặt vào so sánh, thành tích có thể tăng một thời gian nhưng lòng tự trọng lại ngày càng phụ thuộc vào việc hơn thua.}{Some children ask for tutoring not to improve but to avoid falling behind peers. When motivation is tied too tightly to comparison, scores may rise for a while while self-worth becomes increasingly dependent on winning.}

\latcat{5. Bài Làm Đẹp Nhưng Mặt Người Xấu Đi}{Em làm bài ra kết quả tốt hơn trước, nhưng mặt mũi lúc nào cũng bí bách, cáu bẳn và mất ngủ. Nếu chỉ nhìn điểm số, người lớn dễ gọi đó là tiến bộ. Nếu nhìn cả con người, phải thừa nhận rằng có những kiểu tiến bộ đang ăn vào sức sống.}{The child's scores improve, yet the face grows tense, irritable, and sleep-deprived. If adults look only at scores, they call it progress. If they look at the whole person, they must admit that some forms of progress eat away at vitality.}

\latcat{6. Một Lớp Nữa Để Chữa Một Nỗi Lo Nữa}{Mỗi lần lo, gia đình lại thêm một lớp. Mỗi lần chưa yên tâm, lại thêm một thầy. Sau cùng, không ai chắc điều gì đang được giải quyết: lỗ hổng kiến thức của học sinh hay cơn bất an của người lớn.}{Each time worry rises, the family adds another class. Each time reassurance fails, they add another teacher. In the end, no one knows what is truly being addressed: the student's knowledge gaps or the adults' anxiety.}

\latcat{7. Điện Thoại Ban Đêm, Học Thêm Ban Ngày}{Ban ngày lịch học dày, ban đêm lại lén cầm điện thoại để tìm chút cảm giác thoát ra. Sáng hôm sau tiếp tục mệt, tiếp tục học kém hiệu quả, rồi lại được kết luận là cần học thêm nhiều hơn. Một vòng luẩn quẩn rất hiện đại hình thành như thế.}{The day is packed with study, and at night the child secretly clings to the phone to feel some escape. The next morning brings more exhaustion, less effective learning, and the conclusion that even more tutoring is needed. That is how a very modern vicious cycle forms.}

\latcat{8. Lớp Thêm Thành Nơi Giữ Con}{Có nhà cho con đi học thêm không chỉ vì học, mà còn vì không muốn con rảnh quá. Khi lớp học trở thành nơi giữ trẻ có học phí cao, mục tiêu giáo dục bắt đầu lệch đi. Học thêm lúc đó không còn là công cụ hỗ trợ mà thành giải pháp quản lý thời gian thay cho giáo dục nề nếp.}{Some families use tutoring not only for learning but because they do not want the child to be idle. When class becomes paid child storage, the educational goal drifts. Tutoring then stops being support and becomes a time-management substitute for real discipline.}

\latcat{9. Được Khen Chịu Khó, Không Ai Hỏi Có Hạnh Phúc Không}{Em được khen là siêng, là chịu khó, là có tương lai. Nhưng hiếm ai hỏi: em có còn thấy mình là một người sống được không. Một nền giáo dục chỉ biết khen chịu cực mà quên hỏi về sức khỏe tinh thần rất dễ tạo ra những con người giỏi chịu đựng nhưng nghèo niềm vui.}{The child is praised as diligent, enduring, full of promise. Yet almost no one asks whether the child still feels alive. An education system that praises endurance while neglecting mental health easily produces people skilled at suffering but poor in joy.}

\latcat{10. Ít Lớp Hơn Nhưng Sâu Hơn}{Một học sinh cắt bớt một lớp thêm, bớt điện thoại đêm, ngủ đủ hơn và dành thời gian ngồi chữa lỗi sai thật. Sau vài tháng, em không học nhiều nơi hơn, nhưng học chắc hơn. Đây là mẫu truyện ít được tung hô, vì nó không ồn ào bằng lịch kín. Nhưng nhiều khi đó mới là con đường tử tế.}{One student cuts one tutoring class, reduces late-night phone use, sleeps more, and spends time genuinely correcting mistakes. After a few months the child is not studying in more places, but learning more solidly. This story is rarely celebrated because it is not as loud as a packed schedule. Yet it is often the healthier path.}

\begin{goigiaovien}
Những mẫu truyện kiểu này không nhằm kết tội học thêm như một thứ xấu tuyệt đối.
\textit{(These story patterns are not meant to condemn tutoring as something absolutely bad.)}

Chúng nhằm giúp giáo viên, phụ huynh và học sinh nhìn ra sự khác nhau giữa hỗ trợ thật và sự quá tải được ngụy trang bằng lời khen.
\textit{(They are meant to help teachers, parents, and students distinguish real support from overload disguised as praise.)}
\end{goigiaovien}

\begin{tongketchuong}
Không ít đứa trẻ đang được khen là chăm chỉ trong khi thực chất chỉ đang sống trong một guồng mệt mỏi được tổ chức rất đẹp.
\textit{(Many children are praised as hardworking when in truth they are living inside a beautifully organized machine of exhaustion.)}
\end{tongketchuong}

\ngancach